%\documentclass[12pt,serif]{beamer}
%\documentclass[tikz,12pt,svgnames]{beamer}
\documentclass[table,handout,tikz,12pt,svgnames]{beamer}
\usepackage{CM-preamble}
\subtitle{\Huge Les Fichiers en C}
\date{CM5}

%README TODO STUFF

\begin{document}

\begin{frame}
	\titlepage
\end{frame}


%\begin{frame}[fragile=singleslide]
%	\frametitle{Les fichiers en C}
%		\begin{block}{Pas de fichiers dans le langage C}
%			\begin{itemize}
%				\item Bibliothèque standard \texttt{libc.so/libc.a} :\\inclure le fichier entête \texttt{<stdio.h>}
%			\end{itemize}
%		\end{block}
%		\begin{block}{Fichier C = une suite d'octets (= flot)}
%			\begin{itemize}
%			\item Pas de fichiers typés
%			\item Fichiers texte
%			\begin{itemize}
%				\item Les octets représentent des caractères codant les données\footnotesize(souvent le très limité ASCII, mais aussi le populaire et recommandé \textbf{UTF-8})
%				\small
%				\item Standard -- Éditables -- Imprimables
%			\end{itemize}
%			\end{itemize}
%		\end{block}
%\end{frame}

\begin{frame}[fragile=singleslide]
	\frametitle{Les fichiers en C}
		\begin{block}{Pas de fichiers de base dans le langage C}
			\begin{itemize}
				\item Mais dans la bibliothèque standard \texttt{libc.so/libc.a}\\
                    en incluant le fichier d'en-tête \texttt{<stdio.h>}
			\end{itemize}
		\end{block}
		\begin{block}{Un fichier C = une suite d'octets (= flot)}
			\begin{itemize}
                \item Pas de types de fichiers (e.g. image, tableau \ldots)
%			\item Fichiers texte
			\end{itemize}
		\end{block}
		\begin{block}{Fichiers texte}
			\begin{itemize}
				\item Les octets représentent des caractères codant les données \footnotesize(souvent le très limité ASCII, mais aussi le populaire et recommandé \textbf{UTF-8})
				\small
				\item Standard -- Éditables -- Imprimables
%			\end{itemize}
			\end{itemize}
		\end{block}
\end{frame}


\begin{frame}[fragile=singleslide]
	\frametitle{Fichiers Binaires}
%	\begin{block}{Fichiers binaires}
	\vspace{-0.75cm}
		\begin{itemize}
			\item Les octets représentent la copie exacte des données en mémoire sur un système donné
			\item Non standard -- Non éditables -- Non imprimables
			\item Mais lecture / écriture plus rapides (pas d'analyse)
			\item En général, plus compacts
			\begin{itemize}
				\item \small ex: 654875 = 6 octets (\texttt{char}), 2/4 octets (\texttt{short}/\texttt{int})
			\end{itemize}
		\end{itemize}
%	\end{block}
	\begin{block}{Pas d'attribut \textit{texte} ou \textit{binaire} sur un fichier \\(\small dépend de l'interprétation des octets)}
		\begin{itemize}
			\item N'intervient pas à la déclaration
			\item Lié aux opérations applicables
		\end{itemize}
	\end{block}
\end{frame}


\begin{frame}[fragile=singleslide]
	\frametitle{Le type \texttt{FILE}}
	\begin{block}{}
		{\includegraphics[scale=0.9]{../common-images/file1.pdf}}		
		\begin{itemize}
			\item Défini dans \texttt{<stdio.h>}
			\item Structure C contenant
			\begin{itemize}
				\item Identification du fichier associé (descripteur)
				\item Position du curseur dans le fichier
				\item Tampon de lecture / écriture
				\item Indication de mode d'ouverture \ldots
			\end{itemize}
			\item Opérations sont effectuées sur un \texttt{FILE *} fourni à l'ouverture
		\end{itemize}
	\end{block}
\end{frame}


\begin{frame}[fragile=singleslide]
	\frametitle{Fichiers texte: ouverture}
	\begin{block}{}
		\begin{itemize}
			\item Défini dans \texttt{<stdio.h>}
			\item \mintinline{C}{FILE * fopen(char *nom, char *mode)} où
%			\item \texttt{FILE * fopen(char *nom, char *mode)} où
%			\\où \texttt{mode} = \texttt{r}: lecture, \texttt{w}: création, \texttt{a}: allongement
			\vspace{-0.5em}
			\small
			\[
			\hspace{-4cm}
			\texttt{mode}= 
			\begin{cases}
			\texttt{'r'}: & \text{lecture}\\
			\texttt{'w'}: & \text{création/écriture}\\
            \texttt{'a'}: & \text{allongement (ajout à la fin du fichier)}
			\end{cases}
			\]
			\vspace{-1em}
		\end{itemize}
	\end{block}
	{\includegraphics[scale=0.6]{../common-images/file2.pdf}}
\end{frame}


\begin{frame}[fragile=singleslide]
	\frametitle{Fichiers texte: ouverture/fermeture}
	\begin{block}{Retour}
		\begin{itemize}
			\item \texttt{FILE * } si tout va bien
			\item \texttt{NULL} si erreur (ex: fichier inexistant, pas les droits de lecture ou écriture, \ldots)
		\end{itemize}
	\end{block}
	\begin{block}{Fermeture}
		\begin{itemize}
			\item \texttt{fclose(FILE *fp)}
			\item Déconnecte fp du fichier physique
			\item Libère la mémoire du programme associée au fichier
			\item Permet aux autres applications d'utiliser le fichier
		\end{itemize}
	\end{block}
\end{frame}

%NATHALIE: tester pointeur avant fermeture! Sinon, segfault
\begin{frame}[fragile=singleslide]
	\frametitle{Fichiers texte: exemple}
		\begin{block}{} %Gestion dynamique de la mémoire}
			\vspace{-0.5cm}
			\inputminted[
			%		frame=lines,
			%		framesep=2.5mm,
			baselinestretch=0.7,
			fontsize=\footnotesize,
			linenos,
			tabsize=2,
			%		bgcolor=light-gray
			]
			{c}{../codes/file1.c}
		\end{block}
		\footnotesize
		\begin{block}{} %Gestion dynamique de la mémoire}
				\begin{columns}[b]
%					\hspace{-0.3cm}
					\begin{column}{0.7\textwidth}
						\underline{\textit{si toto.txt existe}}\\
						\texttt{toto.txt accessible:fp=0x1d12010\\Fermer toto}
					\end{column}
					\hspace{-2cm}
					\vrule{}
					\hspace{0.5cm}	
					\begin{column}{0.4\textwidth}
						\underline{\textit{si toto.txt n'existe pas}}\\						
						\texttt{toto.txt inaccessible: fp=NULL}
					\end{column}
				\end{columns}
		\end{block}					
\end{frame}


\begin{frame}[fragile=singleslide]
	\frametitle{Fichiers texte: utilisation}
	\begin{block}{}
		\begin{itemize}
			\item Généralisation des manipulations effectuées sur l'entrée/sortie standard (\texttt{stdin}, \texttt{stdout})
			\item Dans \texttt{<stdio.h>}
			\begin{itemize}
				\item entrée standard : \texttt{FILE * stdin}
				\item sortie standard : \texttt{FILE * stdout}
			\end{itemize}
        \item Connexion à l'exécution aux entrées / sorties standard fournies par le système (console par défaut, redirigeables par \texttt{<}, \texttt{>} ou \texttt{|})
			\item Lectures et écritures à partir de la position suivant le curseur
		\end{itemize}
	\end{block}
\end{frame}


%NATHALIE: non null for feof is incorrect, it's non zéro
\begin{frame}[fragile=singleslide]
	\frametitle{Fichiers texte: lecture}
	\vspace{-0.8cm}
	\begin{block}{}
		\begin{itemize}

			\item \texttt{char getc (FILE *fp)} \\
			\hspace{1cm} \footnotesize{\texttt{getchar()} $\Leftrightarrow$ \texttt{getc(stdin)}}
%			\begin{itemize}
%				\item \texttt{getchar()} $\Leftrightarrow$ \texttt{getc(stdin)}
%			\end{itemize}

%			\item \texttt{int fscanf(FILE *fp, char *format, \ldots)}
%			\begin{itemize}
%				\item scanf(\ldots) $\Leftrightarrow$ fscanf(stdin, \ldots)
%				\item Retourne le nombre d'items lus, à vérifier à \underline{chaque utilisation}
%			\end{itemize}
			\item \normalsize \texttt{\textbf{int fscanf}(FILE *fp, char *format, \ldots)} \\
			\hspace{1cm}\footnotesize \texttt{scanf(\ldots)} $\Leftrightarrow$ \texttt{fscanf(stdin, \ldots)} \\
			\footnotesize Retourne le nombre d'items lus, \textbf{à vérifier à \underline{chaque utilisation}}

			
			\item \small \texttt{char * fgets(char *chaine, int taille,FILE *fp)} \\
%			\hspace{7em} \texttt{FILE *fp)}
			\item \texttt{int feof(FILE *fp)} \\
%			\begin{itemize}
				%\item 
				\footnotesize Retourne une valeur différent à zéro si la fin du fichier a été rencontrée lors d'une opération de lecture
%			\end{itemize}
			\item \texttt{int ferror(FILE *fp)} \\
%			\begin{itemize}
%				\item 
				\footnotesize Retourne une valeur différent à zéro si un erreur a été rencontrée lors d'une opération de lecture
%			\end{itemize}
			 \item \texttt{void clearerr(FILE *stream)} \\ remet à zéro l'indicateur d'erreur pour le flux de caractères considéré
		\end{itemize}
	\end{block}
\end{frame}


\begin{frame}[fragile=singleslide]
	\frametitle{Fichiers texte: mode écriture/allongement}
	\begin{block}{}
		\begin{itemize}
			\item \texttt{int putc(char c, FILE *fp)}
			\begin{itemize}
				\item \texttt{putchar(c)} $\Leftrightarrow$ \texttt{putc(c, stdout)}
			\end{itemize}
			\item \texttt{int fprintf(FILE *fp, char *format, \ldots)}
			\item \texttt{int fputs(char *chaine, FILE *fp)}
		\end{itemize}
	\end{block}
\end{frame}


%NATHALIE: tester pointeur avant fermeture! Sinon, segfault
\begin{frame}[fragile=singleslide]
	\frametitle{Fichiers texte: exemple (MAUVAISE)}
	\begin{block}{} %Gestion dynamique de la mémoire}
		\vspace{-0.5cm}
		\inputminted[
		%		frame=lines,
		%		framesep=2.5mm,
		baselinestretch=0.7,
		fontsize=\footnotesize,
		linenos,
		tabsize=4,
		%		bgcolor=light-gray
		]
		{c}{../codes/file3.bad.c}
	\end{block}	
\end{frame}

\begin{frame}[fragile=singleslide]
	\frametitle{Fichiers texte: exemple (BONNE)}
	\begin{block}{} %Gestion dynamique de la mémoire}
		\vspace{-0.5cm}
		\inputminted[
		%		frame=lines,
		%		framesep=2.5mm,
		baselinestretch=0.7,
		fontsize=\footnotesize,
		linenos,
		tabsize=4,
		%		bgcolor=light-gray
		]
		{c}{../codes/file3.good.c}
	\end{block}	
\end{frame}

\begin{frame}[fragile=singleslide]
	\frametitle{Fichiers texte: exemple (clauses de garde)}
	\begin{block}{} %Gestion dynamique de la mémoire}
		\vspace{-0.7cm}
		\inputminted[
		%		frame=lines,
		%		framesep=2.5mm,
		baselinestretch=0.1,
		fontsize=\footnotesize,
		linenos,
		tabsize=2,
		%		bgcolor=light-gray
		]
		{c}{../codes/file3.good2_garde.c}
	\end{block}	
\end{frame}



\begin{frame}[fragile=singleslide]
	\frametitle{Fichiers binaires: ouverture}
%	\begin{block}{}
		\begin{itemize}
			\item Octets représentent la copie exacte du codage des données en mémoire
		\end{itemize}
	\begin{block}{Ouverture}
		\begin{itemize}
			\item \mintinline{c}{FILE *fopen (char *nom, char *mode)} où
			\vspace{-0.5em}
			\small
			\[
			\hspace{-4cm}
			\texttt{mode}= 
			\begin{cases}
			\texttt{rb}: & \text{lecture}\\
			\texttt{wb}: & \text{création/écriture}\\
			\texttt{ab}: & \text{allongement}\\
			\texttt{rb+}: & \text{lecture/écriture}	\\\ldots
			\end{cases}
			\]
			\vspace{-1em}
		\end{itemize}
	\end{block}
	{\includegraphics[scale=0.48]{../common-images/file3.pdf}}
\end{frame}

%NATHALIE: cast de trop dans fwrite & sizeof(<T>)
\begin{frame}[fragile=singleslide]
	\frametitle{Fichiers binaires: fermeture/écriture}
	\begin{block}{Fermeture \footnotesize(idem fichiers texte)}
		\begin{itemize}
			\item \texttt{fclose(FILE *fp)}
			\item \texttt{int feof(FILE *fp)}
		\end{itemize}
	\end{block}
	\begin{block}{Écriture (mode création ou lecture/écriture)}
		\begin{itemize}
%			\item \texttt{int fwrite(void *pt, int taille, int nb,}\\
%			\hspace{6em} \texttt{FILE *fp)}\\
			\item \texttt{size\_t fwrite(const void *pt, \\ \hspace{7em} size\_t taille, size\_t nb, \\ \hspace{7em} FILE *fp);}\\
			Écrit sur le fichier \texttt{fp}, à partir de la position suivant le curseur, \texttt{nb} objets, chacun de taille \texttt{taille}, qui se trouvent contiguëment dans la zone mémoire pointée par \texttt{pt}. \underline{Et la valeur de retour ?}
			\item Utilisation courante :\\
			\texttt{FILE *fp; <T> x; \\fwrite(\&x, sizeof(x), 1, fp);}
		\end{itemize}
	\end{block}
\end{frame}


\begin{frame}[fragile=singleslide]
	\frametitle{Fichiers binaires: écriture}
	\begin{block}{Mode lecture/écriture}
		\begin{itemize}
			\vspace{2cm}
%			\item %\huge TODO IMAGE\\
%			\normalsize
			\item \texttt{char x = '?';}\\
			\texttt{fwrite(\&x, 1, 1, fp);}
			\vspace{3cm}			
%			\item \huge TODO IMAGE
		\end{itemize}
	\end{block}
\end{frame}


\begin{frame}[fragile=singleslide]
	\frametitle{Fichiers binaires: lecture}
	\begin{block}{}
		\begin{itemize}
%			\item \texttt{int fread(void *pt, int taille, int nb, FILE *fp)}\\
%			Lire nb objets de taille <taille> et les copier dans l'espace pointé par pt\\
			\item \texttt{size\_t fread(void *pt, size\_t taille, \\ \hspace{6.5em} size\_t nb, FILE *fp);}\\
			Lire \texttt{nb} objets de taille \texttt{taille} et les copier dans l'espace pointé par \texttt{pt}.\\
			\vspace{0.5cm}
			\item Utilisation courante :\\
			\texttt{<T> x; \\fread(\&x, sizeof(<T>), 1, fp);}
		\end{itemize}
	\end{block}
\end{frame}


\begin{frame}[fragile=singleslide]
	\frametitle{Fichiers binaire: lecture (MAUVAISE)}
	\begin{block}{} %Gestion dynamique de la mémoire}
		\vspace{-0.5cm}
		\inputminted[
		%		frame=lines,
		%		framesep=2.5mm,
		baselinestretch=0.7,
		fontsize=\footnotesize,
		linenos,
		tabsize=4,
		%		bgcolor=light-gray
		]
		{c}{../codes/file_binary_persons.bad.c}
	\end{block}	
\end{frame}
\begin{frame}[fragile=singleslide]
	\frametitle{Fichiers binaire: lecture (BONNE)}
%	\begin{block}{} %Gestion dynamique de la mémoire}
%		\vspace{-0.5cm}
%		\inputminted[
%		%		frame=lines,
%		%		framesep=2.5mm,
%		baselinestretch=0.7,
%		fontsize=\footnotesize,
%		linenos,
%		tabsize=4,
%		%		bgcolor=light-gray
%		]
%		{c}{../codes/file_binary_persons.good.c}
%	\end{block}	
\end{frame}



\begin{frame}[fragile=singleslide]
	\frametitle{Fichiers binaires: accès direct}
	\begin{block}{}
		\begin{itemize}
			\item \mintinline{c}{int fseek(FILE *fp, long déplacement,}\\
			\hspace{11em}\mintinline{c}{int origine)}
%			\normalsize où
%			\vspace{-0.5em}
			\small
			\[
			\hspace{-0cm}
			\text{où   }
			\hspace{0.5em}
			\texttt{    origine}= 
			\begin{cases}
			\texttt{SEEK\_SET}: & \text{début}\\
			\texttt{SEEK\_CUR}: & \text{position courant}\\
			\texttt{SEEK\_END}: & \text{fin}
			\end{cases}
			\]
%			\vspace{-1em}
			\item Positionne le curseur pour la prochaine lecture ou écriture
			\item Position = déplacement + origine
			\item Usage courant : \\
			\texttt{fseek(fp, i*sizeof(<T>), SEEK\_SET);}
            \item \texttt{stdin} et \texttt{stdout} ne supportent pas \texttt{fseek}
		\end{itemize}
	\end{block}
\end{frame}


\begin{frame}[fragile=singleslide]
	\frametitle{Fichiers: conclusion}
	\begin{block}{}
		\begin{itemize}
			\item \textit{Texte} ou \textit{binaire} n'est pas un attribut de fichier
			\item Un fichier texte peut être exploité en binaire comme simple suite d'octets
			\begin{itemize}
				\item ex : \texttt{fseek(fp, i*sizeof(char), SEEK\_SET);}
				\item ex : utilisation de \texttt{fread} ou \texttt{fwrite} sur un fichier texte
			\end{itemize}
			\item Exploitation d'un fichier binaire en texte ?????
		\end{itemize}
	\end{block}
\end{frame}


%mode	description	starts..
%r	rb		open for reading (The file must exist)	beginning
%w	wb		open for writing (creates file if it doesn't exist). Deletes content and overwrites the file.	beginning
%a	ab		open for appending (creates file if it doesn't exist)	end
%r+	rb+	r+b	open for reading and writing (The file must exist)	beginning
%w+	wb+	w+b	open for reading and writing. If file exists deletes content and overwrites the file, otherwise creates an empty new file	beginning
%a+	ab+	a+b	open for reading and writing (append if file exists)	end


%TODO: ANNEXE CON TODOS LOS TIPOS DE APERTURA

%\begin{frame}[fragile=singleslide]
%\end{frame}
%\begin{frame}[fragile=singleslide]
%\end{frame}




% % % % % % % % % % % % % % % % % % % % % % % % % % %
% END
% % % % % % % % % % % % % % % % % % % % % % % % % % %
\end{document}
% % % % % % % % % % % % % % % % % % % % % % % % % % %
% END
% % % % % % % % % % % % % % % % % % % % % % % % % % %
